\section{Результаты}

\subsection{Сборка и проверка}

Программа проверена в Ubuntu 22.04 под WSL2 на процессоре AMD Ryzen 5 7535HS
(6 ядер, 12 логических процессоров, кэш L3 16 МБ, 8 ГБ памяти, выделенных WSL).
Команды выполняются из корня репозитория:
\begin{verbatim}
make
make test
\end{verbatim}
В журнале \texttt{logs/tests.txt}: 12 сценариев, 46 успешных проверок.
Проверяются решение с перестановкой строк, совпадение ответа при разном числе
потоков, вырожденные и несовместные системы, ошибки формата, ключей и записи
(\texttt{/dev/full}), TID потоков, число потоков в \texttt{/proc}, отказ
\texttt{pthread\_create} при \texttt{ulimit -u 1} и ручной ввод с терминала
(через псевдотерминал \texttt{script}).

\lstinputlisting[language={},title={Окончание журнала тестов}, firstline=34]{inc/tests.txt}

\subsection{Демонстрационный сеанс}

\lstinputlisting[language={},title={Сеанс: пример 4 $\times$ 4, случайная система 2000 $\times$ 2000, вырожденная система}]{inc/demo.txt}

Строки, начинающиеся с \texttt{\$}, добавлены при оформлении: это команды
запусков из \texttt{logs/demo.txt}. Решение примера --- $x = (55, 68, 85, 122)/33$.
Поток 0 имеет TID, равный PID: главный поток участвует в вычислениях.

\subsection{Число потоков средствами ОС}

\lstinputlisting[language={},title={Потоки работающей программы при \texttt{-t 4} (\texttt{tests/threads\_demo.sh})}]{inc/threads.txt}

\texttt{ps -L} выводит по строке на поток: столбец LWP --- TID, NLWP --- число
потоков процесса, PSR --- номер процессора, на котором поток выполнялся
последним. Все четыре потока работают на разных процессорах. Поле
\texttt{Threads} в \texttt{/proc/<pid>/status} и каталог \texttt{/proc/<pid>/task}
показывают те же 4 потока, \texttt{top -H} --- загрузку каждого.
Тест 10 опрашивает \texttt{/proc/<pid>/status} каждые 10 мс: при \texttt{-t 1},
\texttt{-t 4}, \texttt{-t 6} наибольшее число потоков равно значению ключа.

\subsection{Системные вызовы}

\lstinputlisting[language={},title={Трассировка \texttt{strace -f -e trace=clone3,futex,exit ./gauss -t 3 -g 3}}]{inc/strace.txt}

Главный поток (1026) дважды вызывает \texttt{clone3} с флагами
\texttt{CLONE\_VM|CLONE\_THREAD}: новые задачи разделяют адресное пространство
и принадлежат тому же процессу. Первые \texttt{futex} --- стартовые ворота
(условная переменная). Далее на каждом барьере потоки засыпают
(\texttt{FUTEX\_WAIT}) на одном адресе, а последний пришедший будит всех
(\texttt{FUTEX\_WAKE} с числом 2147483647). Ответ \texttt{EAGAIN} означает, что
барьер открылся раньше, чем поток успел заснуть. Система $3 \times 3$ даёт
3 шага по 2 барьера. Завершение потоков --- \texttt{exit}; главный поток ждёт их
в \texttt{pthread\_join} (\texttt{futex} на адресе TID).

\subsection{Исследование ускорения и эффективности}

Ускорение $S_p = T_1 / T_p$, эффективность $E_p = S_p / p$, где $T_p$ --- время
решения в $p$ потоках. $T_1$ измерено той же программой с \texttt{-t 1}.
Каждая точка --- медиана пяти запусков на одной случайной системе
(\texttt{tests/bench.sh}, журнал \texttt{logs/bench.txt}).
Измеряется только решение: генерация и проверка в замер не входят.

\begin{center}
\small
\begin{tabular}{r|r|rrrrrrr}
$n$ & $T_1$, мс & \multicolumn{7}{c}{ускорение $S_p$ при $p$ потоках} \\
 & & 2 & 3 & 4 & 6 & 8 & 12 & 16 \\
\hline
100  & 0,21  & 0,04 & 0,03 & 0,02 & 0,01 & 0,01 & 0,01 & 0,005 \\
200  & 0,99  & 0,09 & 0,06 & 0,04 & 0,03 & 0,02 & 0,02 & 0,01 \\
500  & 12,4  & 0,36 & 0,27 & 0,21 & 0,14 & 0,11 & 0,08 & 0,06 \\
1000 & 96,2  & 0,90 & 0,79 & 0,67 & 0,51 & 0,39 & 0,24 & 0,21 \\
2000 & 1705  & 1,53 & 1,84 & \textbf{1,91} & 1,77 & 1,65 & 1,49 & 1,25 \\
3000 & 7151  & 1,42 & 1,91 & \textbf{2,24} & 1,79 & 1,95 & 1,83 & 1,84 \\
\end{tabular}
\end{center}

\pgfplotsset{
  table/comment chars={\#},
  % keeps only the rows with the given n
  only n/.style={x filter/.append code={%
    \edef\rowvalue{\thisrow{n}}\edef\wanted{#1}%
    \ifx\rowvalue\wanted\else\def\pgfmathresult{nan}\fi}},
  unbounded coords=discard,
  bench/.style={width=0.7\textwidth, height=6.5cm, grid=major,
    xlabel={число потоков $p$}, xtick={1,2,3,4,6,8,12,16},
    legend pos=outer north east, legend cell align=left, mark size=1.5pt},
}

\begin{figure}[H]
\centering
\begin{tikzpicture}
\begin{axis}[bench, ylabel={ускорение $S_p$}, ymin=0, ymax=3]
\addplot[black, dashed, domain=1:3, samples=2] {x};
\addlegendentry{идеальное}
\addplot table[x=threads, y=speedup, only n=500]  {inc/bench.txt}; \addlegendentry{$n=500$}
\addplot table[x=threads, y=speedup, only n=1000] {inc/bench.txt}; \addlegendentry{$n=1000$}
\addplot table[x=threads, y=speedup, only n=2000] {inc/bench.txt}; \addlegendentry{$n=2000$}
\addplot table[x=threads, y=speedup, only n=3000] {inc/bench.txt}; \addlegendentry{$n=3000$}
\end{axis}
\end{tikzpicture}
\caption{Ускорение в зависимости от числа потоков и размера системы}
\end{figure}

\begin{figure}[H]
\centering
\begin{tikzpicture}
\begin{axis}[bench, ylabel={эффективность $E_p$}, ymin=0, ymax=1.05]
\addplot table[x=threads, y=efficiency, only n=500]  {inc/bench.txt}; \addlegendentry{$n=500$}
\addplot table[x=threads, y=efficiency, only n=1000] {inc/bench.txt}; \addlegendentry{$n=1000$}
\addplot table[x=threads, y=efficiency, only n=2000] {inc/bench.txt}; \addlegendentry{$n=2000$}
\addplot table[x=threads, y=efficiency, only n=3000] {inc/bench.txt}; \addlegendentry{$n=3000$}
\end{axis}
\end{tikzpicture}
\caption{Эффективность в зависимости от числа потоков и размера системы}
\end{figure}

\subsection{Объяснение результатов}

{\bfseries Цена синхронизации.} На каждом шаге два барьера, всего $2n$.
У \texttt{pthread\_barrier\_t} ожидающий поток засыпает в ядре, а последний
пришедший будит остальных; каждое пробуждение --- системный вызов и
планирование потока, в WSL это десятки микросекунд. При $n = 100$ решение в
одном потоке занимает 0,21 мс, а в двух --- 5,5 мс: 200 барьеров стоят около
5 мс. С ростом $p$ надо будить больше потоков, и время растёт примерно
линейно (при $n = 100$: 5,5; 11,6; 22,5 мс для 2, 4, 8 потоков).

{\bfseries Малые и средние системы ($n \le 1000$).} Работа шага равна
примерно $(n-k)^2$ операций и делится на $p$, а затраты на барьеры от $n$
почти не зависят. Пока шаг короче пробуждения потоков, параллельная версия
медленнее последовательной; при $n = 1000$ два потока уже почти сравнялись
с одним ($S = 0{,}90$).

{\bfseries Большие системы ($n = 2000$--$3000$).} Работа шага больше затрат
на синхронизацию, и потоки дают выигрыш: до 1,9 при $n = 2000$ и 2,2 при
$n = 3000$ на 4 потоках. Дальше рост ограничивают две причины. Матрица
(32 и 72 МБ) не помещается в кэш L3, и каждый шаг заново читает и пишет
оставшуюся часть матрицы из оперативной памяти; пропускная способность памяти
общая для всех ядер. Это видно уже в одном потоке: при росте $n$ с 1000 до 2000
работа растёт в 8 раз, а время --- в 18. Кроме того, последние шаги почти не
содержат работы, а барьеры остаются.

{\bfseries Число потоков больше 4--6.} У процессора 6 физических ядер; два
логических процессора одного ядра делят его вычислительные блоки и кэши.
При $p > 6$ вычислений не прибавляется, а барьеров на поток больше; при
$p = 16$ потоков больше, чем процессоров, и потоки ждут ещё и очереди на
процессор. Эффективность монотонно падает с ростом $p$ при любом $n$.

{\bfseries Последовательная часть.} Обратный ход, поиск масштаба матрицы и
выбор ведущей строки ($O(n^2)$ в сумме) выполняются одним потоком, но при
$n = 3000$ занимают миллисекунды из секунд. Значения элементов матрицы на время
почти не влияют: метод выполняет одинаковое число операций для любой
невырожденной плотной матрицы.

Разброс отдельных замеров в виртуальной машине WSL достигает десятков процентов
(например, $n = 3000$, $p = 6$ выбивается из ряда).
